\documentclass{article}
\usepackage{amsmath}
\begin{document}
\section*{Actuation slew and independent evidence of the owed heading}
\paragraph{Material change.}
Keep the explicitly modeled forward-only wheel dynamics of PR 846,
with sample period one, radius one, spacing two, initial speeds four,
and both wheels at rest after step five. A locally authorized wheel
may saturate excessive braking to its actual bound; the other must
follow an admissible target. Optional limits on emitted linear and angular
target differences are separate from physical wheel saturation.
The two fault states have weak left or weak right respectively,
with brake bounds \((1,2)\) and \((2,1)\).
Initially study the earlier correlated obligations (straight in the
first state, turn by one radian in the second). Then permit an
authorized dispatcher to choose either heading in either state.
Neither that choice nor private fault identity is baked into
the actuator command or observation. No motion capability or
physical feedback at the robot is inferred from a commanded target.

\paragraph{Actual configured command path.}
An independent one-unit maximum step in body forward and angular
commands is equivalent to a two-unit \(\ell_1\) step bound on
the desired wheel pair, measured against the previous emitted
targets. Under the actual pinned Gazebo selected source with this
configuration, a previously viable unconstrained blind request for
\((2,1)\) after \((3,3)\) instead writes \((2.5,1.5)\).
The model does not substitute either target for measured wheel speed.
Exact integer paired-state search and independent coordinate-word
enumeration agree. Under the unit scalar configuration the correlated
zero/one-radian task is infeasible on the integer grid with neither
or left-only saturation, needs arc 11 with only right saturation,
and needs arc 10 with both. The unconfigured both-port optimum is
8 from PR 846. Native source writes verify the two positive
bounded profiles. The 10 and 11 lower bounds are scoped to the
integer request grid; the native command type permits nonintegers.

\paragraph{Ordering obstruction even for real requests.}
Under any common admissible request word, the stronger left wheel
cannot end a step faster than the weaker left wheel, and the weak
right cannot end slower than the strong right. This follows inductively
from monotonicity of \(v\mapsto \max(u,v-b)\) and its reversed
dependence on \(b\); no-clipping when declared simply restricts
which common \(u\)'s are legal. Distances preserve both inequalities.
In the weak-right state the heading change is therefore always at
least the weak-left state's heading change. Reversing the correlated
obligations (turn in weak left, straight in weak right) is not
realizable by a blind common word, even without the scalar limit
or distance budget. In the earlier correlation the branch could help.
A fault label without the direction of the owed effect is insufficient.

\paragraph{Independent authority evidence is not physical feedback.}
Hold the wheel fault fixed and vary only the authorized heading.
Every legal physical observation can be identical, even for an
adaptive controller, because the plant transition is independent
of the unobserved authority choice. Induction on the common
observations gives the same issued commands (or same law of commands
with common randomness) and the same heading, so no physical-only
observer guarantees both. Even perfect wheel identification does
not identify who owes the turn. This lower bound assumes no
physical signal of the lawfully chosen goal; supplying one is an
interface change.

Conversely, one trustworthy heading bit, without any wheel-state
feedback, is enough by step two at arc nine. A common first two
steps are \((3,3),(2,2)\). Thereafter the straight branch chooses
\((1,1),(0,0),(0,0)\), and the turn branch chooses
\((1,2),(0,1),(0,0)\). Both branches track exactly across both
faults, including the scalar command bounds. Arc nine is necessary
even with advance goal and fault knowledge for the turned weak-left
case: its left distance is at least eight and the required
right-left excess is two.

\paragraph{Why step three is too late for the heading-only interface.}
For a fixed revealed heading, a blind-to-wheel branch must realize
the same final heading in both hardware states. The directional
inequalities above force all realized wheel speeds to match
throughout that entire branch, since the nonnegative accumulated
differences must sum to zero. Common realized braking is at most
one. If the first usable heading bit is at step three, sharing the
prefix leaves two intervals. Let \(p,q\) be contributions to
doubled distances excluding step-four speeds, and \((x_i,y_i)\)
be those speeds for the goal branches. The two heading
constraints give
\[
q+2y_0=p+2x_0,\qquad q+2y_1=p+2x_1+4.
\]
Since \(x_0,y_1\le1\) to stop next tick and all speeds are
nonnegative, equality forces \(x_1=y_0=0\) and \(q-p=2\).
Reaching those zeros from the common step-three state requires
both its wheel speeds to be at most one. From four under three
unit-brake steps both must equal one with the identical fastest
prefix, so \(p=q\), contradiction. This bound also holds for
real sampled velocities, even without the command slew setting.
It does not cover controllers that both see the goal and use
informative physical feedback to choose different later words
by hardware state.

\paragraph{Checks and residuals.}
Twelve finite settings and 479 feasible heading-pair queries
match independently. Eight witness profiles produce 40 actual
desired wheel writes under the configured native scalar limiter;
the local saturated response is separate and stipulated. A
joint controller with the same plant and observation law ties.
No authenticated delivery of the heading bit is provided here,
and native target writes do not prove tracking or clearance.
This is a scoped R6 and authority-interface result, not hardware,
independent mission, or I1 acceptance.
\end{document}
